\section{Leakage audit protocol}\label{app:leakage}

Look-ahead bias is a common reason why published backtests fail out of sample. This paper's predecessor (an
undergraduate thesis) contained two forms of it, which we document here because they are generic.
First, a ``residual correction'' added a fraction of the realised test error to the forecast,
$\hat P_t\leftarrow\hat P_t+0.6\cdot\mathrm{MA}_3(P_t-\hat P_t)$, which uses $P_t$ to forecast $P_t$.
Second, a multi-day target $\bar r_{t+1:t+3}$ was included among the input features, so the last input
row overlapped two of the three target days. Both produced large apparent gains that vanish without
the leak. Neither survives in the present code base.

\paragraph{How large are the effects?} \Cref{tab:leak} reproduces both mechanisms on clean data. With
the lagged three-day target among the inputs, a gradient-boosted model reaches a directional hit rate of
$88.5\%$ for the next-day return and a cost-free long-only Sharpe ratio of $8.7$; without it the same
model has a hit rate of $52.7\%$ and a Sharpe ratio of $0.43$. The residual correction lowers the RMSE
of a 60-lag price model from $52.6$ to $44.9$ index points ($11\%$ \emph{below} the random walk,
which \cref{prop:bound} shows to be impossible without future information) and lifts its directional
accuracy from $49.8\%$ to $65.1\%$. Both numbers are close to the headline results of the
undergraduate thesis, which indicates that these two mechanisms produced them. More generally, a
forecaster that beats the random walk on price levels by more than a fraction of a percent should be
checked for leakage before anyone uses its results.

\begin{table}[!htbp]\centering\scriptsize
\caption{Anatomy of a leak: the two look-ahead mechanisms of the undergraduate thesis reproduced on CSI 300 (out of sample, 2012--June 2024)}\label{tab:leak}
\begin{adjustbox}{max width=\textwidth}
\begin{tabular}{llrrrrrrr}
\toprule
Mechanism & Target & N & Hit rate & PT $p$ & OOS $R^2$ & Long-only Sharpe & RMSE & DA$^{\Delta}$\\
\midrule
leak-free features & 3-day mean & 3,029 & 0.5239 & 0.0065 & -0.0705 & 0.5110 & -- & --\\
leak-free features & next day & 3,032 & 0.5270 & 0.0019 & -0.0278 & 0.4307 & -- & --\\
features + lagged 3-day target (leak A) & 3-day mean & 3,029 & 0.7999 & 0.0000 & 0.6128 & 6.2158 & -- & --\\
features + lagged 3-day target (leak A) & next day & 3,030 & 0.8851 & 0.0000 & 0.8425 & 8.6909 & -- & --\\
price model, leak-free & price level & 3,032 & 0.4977 & -- & -- & -- & 52.6446 & 0.4992\\
price model + residual correction (leak B) & price level & 3,032 & 0.6511 & -- & -- & -- & 44.8885 & 0.5582\\
random walk & price level & 3,032 & -- & -- & -- & -- & 50.6798 & 0.4962\\
\bottomrule
\end{tabular}
\end{adjustbox}
\par\smallskip\begin{minipage}{\textwidth}\footnotesize
Upper block: gradient boosting on 10-day windows of the causal features, with and without the 3-day forward mean return (lagged one day, as in the original code) among the inputs. Lower block: OLS on 60 lagged price levels, with and without the residual correction $0.6\cdot\mathrm{MA}_3(P_t-\hat P_t)$. Long-only Sharpe ignores costs.
\end{minipage}
\end{table}

We impose five safeguards.

\begin{enumerate}[leftmargin=1.5em]
\item \textbf{Causal features by construction.} Every feature is a backward-looking rolling statistic.
A unit test perturbs all prices after a random date and asserts that every feature up to that date is
bit-for-bit unchanged.
\item \textbf{Targets never enter \texttt{predict}.} Models receive covariates and $r_t$ only; the target
column is not passed. Recursive models (GARCH, Markov switching) update their filters with $r_t$, which
is known at the close of $t$.
\item \textbf{End-to-end corruption test.} For every model class, the full walk-forward is run twice
(on the original data and on data whose prices after a cut-off are multiplied by 1.3) and all
forecasts up to the cut-off must coincide to relative precision $10^{-8}$. The test caught one
violation during development: a Markov-switching fit drew random starting values from the global
random-number generator, making forecasts non-reproducible; the generator is now seeded.
\item \textbf{Execution lag.} Positions decided at the close of $t$ are filled at the open of $t+1$ and
earn only the intraday return of $t+1$; the overnight return accrues to the previous position.
\item \textbf{Second data sources.} Realised measures at row $t$ use only the 5-minute bars of day $t$.
The VIX enters every series with the close of the last US trading day strictly before the local date,
because the US close falls after the Asian and European closes of the same calendar day. A unit test
multiplies all realised measures after a cut-off, and the VIX from the cut-off date onwards, by three
and checks that no feature up to the cut-off changes.
\end{enumerate}

\paragraph{Look-ahead through pretraining.} None of these tests can detect a different channel: a
pretrained foundation model whose training corpus contains the evaluation period may have stored the
future in its weights. The documented corpora of Chronos and TimesFM~2.5 do not list the series we
use, but they cannot be audited series by series. \Cref{sec:lockbox} therefore tests for this channel
with the release dates of the models.

